% ---------------------------------------------------------------------
\section{欧氏红外有限性与重整化群}\label{sec:ee-euclid}

闵氏空间中，无质量传播子在整个光锥上奇异，红外发散来自钳夹面（\cref{sec:ir-landau}与\cref{sec:ir-pinch}）。欧氏空间简单得多：Wick 转动后 $k_E^2=0$ 当且仅当 $k_E=0$，被积函数只在某些内线动量为零处奇异。这些点确实是 Landau 方程的解（\cref{rem:ir-euclid-corners}），是否发散只取决于其附近的幂次计数。称欧氏外动量非例外，如果它们的任何真部分和都不为零；两点函数在 $q\ne0$ 时自动非例外。$\Pi$ 的 $\overline{\rm MS}$ 重整化定义为（\cref{def:prelim-msbar}，$S_\varepsilon=(4\pi)^\varepsilon\ee^{-\varepsilon\gamma_E}$）
\begin{equation}\label{eq:ee-PiMS}
\hat\Pi\equiv\mu^{2\varepsilon}S_\varepsilon^{-1}\,\Pi_{\rm bare}\big|_{g_0^2=4\pi\as(\mu)\mu^{2\varepsilon}S_\varepsilon^{-1}Z_g^2},\qquad
\Pi_{\overline{\rm MS}}(q^2;\mu,\as(\mu);\varepsilon)\equiv\hat\Pi-Z_\Pi(\as(\mu),\varepsilon),
\end{equation}
$Z_\Pi$ 与 $q$ 无关，各阶系数是 $\varepsilon$ 的纯极点，由有限性确定。$\hat\Pi$ 恰是 QCD$\times$QED 中光子自能在 $e$ 的最低阶乘在 $e^2$ 上的因子，$Z_\Pi$ 是光子场重整化的相应部分；它与 $q$ 无关，因为 $\Pi^{\mu\nu}$ 的超表观发散部分是 $q$ 的二次多项式，横向性只允许 $(q^\mu q^\nu-q^2g^{\mu\nu})\times$常数。

\begin{theorem}[欧氏红外有限性 \tagR]\label{thm:ee-euclid-ir}
（$R_{\mathrm{PT}}$：在重整化微扰论中对红外正规化的部分子态严格。）设可重整的场论只含无质量场，且没有超可重整耦合（每个相互作用项的量纲恰为 $4$）。则在非例外的欧氏外动量处，用 $\overline{\rm MS}$（或带辅助质量的 BPHZL）重整化的 Green 函数在微扰论的每一阶都有限。无质量 QCD 在任意协变规范中满足这些条件；特别地，$\Pi_{\overline{\rm MS}}(q^2)$ 在 $q^2<0$ 处每一阶有限，且与规范参数无关。（本定理只涉及部分子态的 Green 函数。）
\uses{thm:qcd-renormalizability,def:prelim-msbar,def:qcd-dimreg,prop:qcd-xi-independence,prop:qcd-superficial-degree,lem:prelim-scaleless}
\end{theorem}
\begin{proofsketch}
\emph{第 1 步（紫外）。}由\cref{thm:qcd-renormalizability}（\tagR，$R_{\mathrm{PT}}$），$\overline{\rm MS}$ 抵消项是被重整子图外动量的多项式，不在非零动量处产生奇点；减除不在零动量处进行（零动量的 BPHZ 减除在无质量理论中会制造红外发散，故 BPHZL 引入辅助质量\cite{LowensteinZimmermann1975}）。

\emph{第 2 步（奇点的位置）。}设 $S$ 是同时取零动量的一组内线，其余内线 $H$ 的动量不为零；$S$ 上的 $L_S$ 个独立软圈动量按 $\lambda\to0$ 标度。记 $N_g,N_q,N_c$ 为软的胶子、夸克、鬼场线数，$V_3=n_{ggg}+n_{qqg}+n_{ccg}$ 与 $V_4$ 为\emph{全部}连线都软的三、四点顶点数。

\emph{第 3 步（红外度）。}测度给出 $\lambda^{4L_S}$；每条软胶子线与鬼场线给出 $\lambda^{-2}$（协变规范中 $(1-\xi)k_\mu k_\nu/k^4$ 一项标度相同），每条软夸克线给出 $\lambda^{-1}$；全软的三胶子顶点与鬼场顶点（$-gf^{abc}p^\mu$，$p$ 为出射鬼场动量）各带因子 $\lambda$；连接 $S$ 与 $H$ 的顶点的动量因子为 $\order{1}$ 或更小，只会改善收敛。于是当
\begin{equation}\label{eq:ee-irdeg}
d(S)\equiv4L_S-2N_g-N_q-2N_c+n_{ggg}+n_{ccg}>0
\end{equation}
时 $S$ 附近的积分按幂次计数收敛。当 $H$ 连通时（硬岛的情形见第 4 步），每个全软顶点给出软动量的一个独立约束（$S$ 的每个连通分支都与 $H$ 相连），故 $L_S=N_S-V_3-V_4$，$N_S\equiv N_g+N_q+N_c$。设 $A$ 为落在与 $H$ 相连的顶点上的软线端点数（附着数），$A_q$、$A_c$ 为其中的夸克与鬼场端点数，则数端点得 $2N_S=3V_3+4V_4+A$，$2N_q=2n_{qqg}+A_q$，$2N_c=2n_{ccg}+A_c$。用 $-2N_g-N_q-2N_c=-2N_S+N_q$，
\begin{equation}\label{eq:ee-irdeg2}
\begin{split}
d(S)&=2N_S+N_q-4V_3-4V_4+n_{ggg}+n_{ccg}\\
&=(3V_3+4V_4+A)+\big(n_{qqg}+\tfrac12A_q\big)-4V_3-4V_4+n_{ggg}+n_{ccg}=A+\tfrac12A_q .
\end{split}
\end{equation}
% VF-EE-13 (writer re-derived)

\emph{第 4 步（$A\ge2$）。}附着一次的分支经唯一的附着线与其余部分相连，动量守恒使该线动量恒为零：这是零动量线上的蝌蚪图，在无质量维数正规化中是无标度积分而为零（\cref{lem:prelim-scaleless}），也被 $\langle A_\mu\rangle=0$ 禁止。若 $H$ 有不含流顶点的分支（“硬岛”），其外线全软，整体按质量量纲标度为 $\lambda^{4-E_g-\frac32E_q-E_c}$；由\cref{prop:qcd-superficial-degree}（\tagR，$R_{\mathrm{def}}$），这恰是把硬岛的线全部算作软线时\cref{eq:ee-irdeg}给出的标度，故把硬岛并入 $S$，$A$ 只计附着在含流顶点的分支上的端点。因此 $d(S)=A+\frac12A_q\ge2$：一切零动量区域都可积。最小情形是两条硬线之间交换一个软胶子：$\lambda^4\cdot\lambda^{-2}$。

\emph{第 5、6 步（收敛定理与 QCD）。}“一切软子集 $d(S)>0$、一切紫外子图已重整”蕴含积分绝对收敛：这是无质量传播子的幂次计数定理，BPHZL 情形见\cite{LowensteinZimmermann1975,Lowenstein:1975ps}，维数重整化情形见\cite{BreitenlohnerMaison1977a,Breitenlohner:1975hg,Breitenlohner:1976te}；非例外动量处逐点有限的物理层面的一般证明见\cite{PoggioQuinn1976}；无质量理论中 Feynman 振幅红外奇异性的早期系统分析见\cite{Symanzik1973,KinoshitaUkawa1976}。\cref{rem:ir-euclid-corners}中的零动量 Landau 解正是第 2 步的位形。无质量 QCD 的每个相互作用项（$\bar\psi\sla A\psi$，$AA\partial A$，$AAAA$，$\partial\bar c\,Ac$）的红外维数都恰为 $4$（夸克 $\frac32$，胶子与鬼场 $1$，导数 $+1$）。$\Pi^{\mu\nu}$ 是 QCD$\times$QED 中光子这一基本场在 $e^2$ 阶的自能，$q\ne0$ 时非例外；规范不变流的关联函数与规范参数无关（\cref{prop:qcd-xi-independence}(b)，\tagR，$R_{\mathrm{PT}}$）。完整论证见\cite{LowensteinZimmermann1975,Lowenstein:1975ps,Breitenlohner:1975hg,Breitenlohner:1976te}与\cite{PoggioQuinn1976}。
\end{proofsketch}

\paragraph{关于所引定理的适用范围。}\cite{Breitenlohner:1976te}要求每个相互作用项的红外次数不小于 $4$；如上所示，这对协变规范中的无质量 QCD（包括 Faddeev--Popov 鬼场项）成立。不过 Breitenlohner--Maison 只在例子中提到带矢量型费米子的 Yang--Mills 场，没有明确讨论鬼场，而且其存在性定理是关于缓增分布的陈述；非例外动量处的逐点陈述是\cite{PoggioQuinn1976}（对 BPHZL 减除为\cite{LowensteinZimmermann1975}）的结果。对本章唯一需要的流两点函数，分布意义下的陈述已经足够：由\cref{lem:ee-single-scale}，它的每一阶是有限个 $(Q^2)^{-k\varepsilon}$ 与一个常数之和，$\varepsilon$ 展开的各个负幂系数是 $L$ 的多项式；若这些系数与任意试验函数的积分都为零，则它们逐点为零（这正是\cref{thm:ee-R-finite}证明中的多项式论证），所以分布意义下的有限极限就是逐点的有限极限。纯 Yang--Mills 中规范不变算符关联函数的全阶紫外与红外控制见\cite{Frob:2015uqy}。

\begin{example}[两圈图的软子集]\label{ex:ee-soft-sets}
以\cref{fig:ee-oneloop}(d) 的胶子交换图为例。记流顶点为 $J_1$（$q$ 流入）与 $J_2$（$q$ 流出），胶子顶点为 $a$、$b$，四条夸克线为 $q_1=J_1a$、$q_2=aJ_2$、$q_3=J_2b$、$q_4=bJ_1$，胶子线为 $g=ab$。能同时取零动量、而其余线保持一般动量的线集合 $S$ 受两条约束：一个顶点上若只剩一条非软线，动量守恒迫使它也为软；$q\ne0$ 必须经由硬线从 $J_1$ 流到 $J_2$。满足约束的 $S$ 恰有九个。$\{g\}$：$\lambda^4\cdot\lambda^{-2}$，$A=2$，$A_q=0$，$d=2$。单条夸克线 $\{q_i\}$（四个）：$\lambda^4\cdot\lambda^{-1}$，$A=A_q=2$，$d=3$。$\{q_1,q_2,g\}$ 与 $\{q_3,q_4,g\}$：各含一个全软顶点，$L_S=3-1=2$，$\lambda^8\cdot\lambda^{-1-1-2}$，$A=3$，$A_q=2$，$d=4$。$\{q_1,q_3\}$ 与 $\{q_2,q_4\}$：$L_S=2$，$\lambda^8\cdot\lambda^{-2}$，$A=A_q=4$，$d=6$。每一种情形的直接计数都与\cref{eq:ee-irdeg2}的 $d=A+\frac12A_q$ 一致，最小值 $2$ 来自两条硬夸克线之间的单个软胶子。在自能图（\cref{fig:ee-oneloop}(e)）中还出现硬岛：当自能圈两侧的夸克线变软、而圈内的胶子与夸克保持一般动量时，自能子图的外线全软；按第 4 步把它并入 $S$，即整条上方夸克线连同自能圈为软，$L_S=2$，$\lambda^{8}\cdot\lambda^{-1-1-1-2}$，$A=A_q=2$，$d=3$。
\uses{thm:ee-euclid-ir}
\end{example}

\begin{proposition}[Adler 函数的重整化群 \tagR]\label{prop:ee-adler-rg}
（$R_{\mathrm{PT}}$：在重整化微扰论中对红外正规化的部分子态严格。）在无质量 QCD 中，以\cref{eq:ee-PiMS}定义 $\Pi_{\overline{\rm MS}}$，记 $L=\ln(Q^2/\mu^2)$，$Q^2=-q^2>0$。则：
(a)~$\big(\mu^2\partial_{\mu^2}+\beta(\as)\partial_{\as}\big)\Pi_{\overline{\rm MS}}=\gamma_\Pi(\as)$ 与 $q$ 无关（是接触项），$\gamma_\Pi=\frac{R_0}{12\pi^2}[1+d_1a+\order{a^2}]$；
(b)~$\big(\mu^2\partial_{\mu^2}+\beta(\as)\partial_{\as}\big)D=0$，从而 $D(Q^2)=R_0\,d\big(a(Q^2)\big)$，$d(a)=1+d_1a+d_2a^2+\cdots$，$a(Q^2)$ 是 $\mu^2=Q^2$ 处的跑动耦合；
(c)~在 $a^n$ 阶，$\Pi_{\overline{\rm MS}}$ 是 $L$ 的多项式，次数不超过 $\max(1,n)$（$d_1b_0\ne0$ 时恰为此次数）；$D$ 的 $a^n$（$n\ge1$）系数是 $L$ 的次数不超过 $n-1$ 的多项式。% VF-EE-9
（本命题只涉及部分子态的 Green 函数。）
\uses{def:ee-adler,prop:prelim-rg,prop:qcd-rge,lem:ee-single-scale,thm:ee-euclid-ir}
\end{proposition}
\begin{proof}
(a)~记 $F\equiv\Pi_{\overline{\rm MS}}=\hat\Pi-Z_\Pi$。固定 $g_0$ 与 $\varepsilon$ 时 $\Pi_{\rm bare}$ 与 $\mu$ 无关，故 $\mu^2\frac{\dd\hat\Pi}{\dd\mu^2}=\varepsilon\hat\Pi$；$D$ 维中 $\mu^2\frac{\dd\as}{\dd\mu^2}=-\varepsilon\as+\beta(\as)$（\cref{prop:qcd-rge}，\tagR，$R_{\mathrm{PT}}$）。于是
\begin{equation}\label{eq:ee-rge-eps}
\mu^2\frac{\dd F}{\dd\mu^2}-\varepsilon F=\varepsilon\big(Z_\Pi+\as Z_\Pi'\big)-\beta(\as)Z_\Pi'\equiv\hat\gamma(\as,\varepsilon),\qquad Z_\Pi'\equiv\partial_{\as}Z_\Pi .
\end{equation}
右边与 $q$ 无关。左边在 $\varepsilon\to0$ 时有有限极限：由\cref{thm:ee-euclid-ir}，$F$ 在 $q^2<0$ 处每一阶有限，而由\cref{lem:ee-single-scale}，它每一阶是有限个 $(Q^2/\mu^2)^{-k\varepsilon}$ 与一个常数之和，$\varepsilon$ 展开系数是 $L$ 的多项式，对 $\mu^2$ 求导保持有限性。故 $\hat\gamma$ 有有限极限 $\gamma_\Pi$。写 $Z_\Pi=\sum_{j\ge1}z_j(\as)\varepsilon^{-j}$，则 $\hat\gamma=(z_1+\as z_1')+\sum_{j\ge1}\varepsilon^{-j}[z_{j+1}+\as z_{j+1}'-\beta z_j']$；有限性迫使方括号为零（'t~Hooft 一致性条件），且 $\gamma_\Pi=\frac{\dd}{\dd\as}(\as z_1)$。取 $\varepsilon=0$ 即得 (a) 的第一句。由 $D=-12\pi^2\partial_L\Pi_{\overline{\rm MS}}$，(a) 也可写成 $\gamma_\Pi=\frac{D}{12\pi^2}+\beta\partial_{\as}\Pi_{\overline{\rm MS}}$；在 $L=0$ 处取值，$\beta=\order{\as^2}$，即得 $\gamma_\Pi$ 的展开。

(b)~固定 $Q^2$ 时 $\mu^2\partial_{\mu^2}=-\partial_L$，(a) 写成 $(-\partial_L+\beta\partial_{\as})\Pi_{\overline{\rm MS}}=\gamma_\Pi$；对 $L$ 求导，右边消失，而 $D=-12\pi^2\partial_L\Pi_{\overline{\rm MS}}$，故 $(-\partial_L+\beta\partial_{\as})D=0$。$D$ 是 $Q^2/\mu^2$ 与 $\as(\mu)$ 的无量纲、有限、与 $\mu$ 无关的函数，由\cref{prop:prelim-rg}（\tagR，$R_{\mathrm{PT}}$；特征线法）$D(L,\as(\mu))=D(0,\as(Q^2))\equiv R_0\,d(a(Q^2))$；$d(0)=1$ 来自自由夸克（\cref{def:ee-adler}）。

(c)~记 $\mu^2\frac{\dd a}{\dd\mu^2}=-b_0a^2-b_1a^3-b_2a^4-\cdots$，$b_0=\frac{\beta_0}4$，$b_1=\frac{\beta_1}{16}$，$b_2=\frac{\beta_2}{64}$（\cref{tab:front-notation}第 8 行）。固定 $\mu$ 时 $\frac{\dd a(Q^2)}{\dd L}=-b_0a(Q^2)^2-\cdots$，$a(Q^2)|_{L=0}=a$，逐阶求解得
\begin{equation}\label{eq:ee-arun}
a(Q^2)=a-b_0La^2+\big(b_0^2L^2-b_1L\big)a^3+\big(-b_0^3L^3+\tfrac52b_0b_1L^2-b_2L\big)a^4+\order{a^5}.
\end{equation}
代入 $D=R_0d(a(Q^2))$，$D$ 的 $a^n$ 系数是 $L$ 的次数不超过 $n-1$ 的多项式，首项 $R_0d_1(-b_0L)^{n-1}$。积分 $D=-12\pi^2\partial_L\Pi_{\overline{\rm MS}}$ 得 $\Pi_{\overline{\rm MS}}(L)=\Pi_{\overline{\rm MS}}(0)-\frac1{12\pi^2}\int_0^L\dd L'\,D(L')$，$\Pi_{\overline{\rm MS}}(0)=\sum_nc_na^n$ 与 $q$ 无关；故 $a^0$ 阶是 $L$ 的一次式，$a^n$（$n\ge1$）阶次数不超过 $n$，$d_1b_0\ne0$ 时首项 $-\frac{R_0d_1(-b_0)^{n-1}}{12\pi^2n}L^n$ 不为零。
\end{proof}
常数 $c_n$（如 $c_0=\frac53\frac{R_0}{12\pi^2}$）依赖于方案，是接触项，不进入 $D$ 与 $R$；$d_n$ 不受 $\Pi$ 的加法重整化影响，但 $n\ge2$ 时依赖于 $\as$ 的方案。
